\section{Introduction}

Malware behavior analysis at scale requires transparent, low-perturbation
observation of software execution. Analysts need to know which files a program
opens, whether it creates child processes, when it calls \texttt{execve}, and
whether it reaches unusual traps or memory-permission changes. Software tracers
such as strace~\cite{strace}, eBPF~\cite{ebpf}, and QEMU~\cite{bellard-qemu} run
inside the environment being studied, making them detectable by evasive workloads
and introducing perturbation that can alter behavior. A workload that detects the
tracer may change its code path, delay execution, or suppress malicious activity
entirely.

Hardware-assisted tracing provides an alternative observation point below the
operating system and the traced process. ARM ETM and Intel PT record
architectural events with low perturbation, enabling transparent analysis that is
invisible to the target workload. On ARM, these capabilities have been used for
transparent malware tracing~\cite{ning2017ninja,ning2019ninja}, kernel module
tracing~\cite{hart2020esorics}, and firmware fuzzing
feedback~\cite{microafl2022icse}. On RISC-V, Raft demonstrates hardware-assisted
dynamic information flow tracking~\cite{raft2023raid}. However, raw hardware
events are not enough to answer security questions: a program counter needs an
executable and load address; a syscall transition needs a process context;
expected behavior may come from a reference run rather than from hardware.

Existing hardware-assisted systems focus on ARM-specific secure-world
isolation~\cite{ning2017ninja,ning2019ninja}, kernel
protection~\cite{hart2020esorics}, online taint
enforcement~\cite{raft2023raid}, or path-coverage
feedback~\cite{microafl2022icse}. These systems demonstrate the value of hardware
tracing, but they do not address the semantic gap between raw events and
auditable behavior summaries. None connect hardware trace events to
source-labeled behavior provenance on RISC-V, and none reconstruct user-space
Linux behavior with explicit evidence labeling that separates hardware-derived
fields from supporting metadata or reference logs.

Our insight is that the software-observer problem and the hardware-semantic-gap
problem should be solved together. Hardware trace events provide an observation
point outside the target process, but they become useful for malware analysis
only when they are connected to OS-level behavior semantics. \rvmt therefore
labels each reconstructed field with its provenance source, enabling auditable
behavior reconstruction without overstating which facts came directly from
hardware.

We present \rvmt, a hardware-assisted behavior tracing system for RISC-V Linux
that connects a CVA6/Genesys2 hardware trace path with an offline reconstruction
pipeline. \rvmt captures architectural events---retire, trap, syscall
entry/return, privilege transition, branch/jump, context, and marker---below the
OS, then joins them with ELF metadata and runtime process maps while labeling
each reconstructed field with its evidence source. The result is a behavior trace
whose provenance is explicit: the reader can see which parts came from hardware
and which parts came from supporting metadata or reference logs.

Realizing this system requires addressing three technical challenges. First, how
to capture the right architectural events (retire, trap, syscall, privilege)
without software interference on a RISC-V core. Second, how to connect raw
program-counter events to executable identity and process context under PIE, ASLR,
fork/exec, and dynamic loading. Third, how to ensure reconstructed behavior is
auditable by labeling each field with its provenance source---hardware trace, ELF
metadata, runtime map, or reference log.

We make the following contributions:

\begin{itemize}
  \item A CVA6/Genesys2 hardware trace path that captures retire, trap, syscall
  entry/return, privilege transition, branch/jump, context, and marker events
  below the OS for transparent RISC-V Linux observation.
  \item An offline behavior reconstruction pipeline that connects raw trace events
  to executable identity and process context, with each reconstructed field
  explicitly labeled by its evidence source.
  \item A code-attribution mechanism that resolves executable identity under
  PIE/ASLR, dynamic loading, fork/exec, and stripped binaries, attributing
  108/108 key PC events to the target ELF/function level and passing 6/6
  mapping edge cases on controlled workloads.
  \item A bounded anti-analysis case study showing that, for one controlled
  anti-debug-like workload where software tracing is visible or perturbs
  \texttt{ptrace}, \rvmt still reconstructs the behavior from Genesys2 BRAM
  trace records.
  \item A reproducible artifact package with manifests and SHA-256 checks that
  enables independent verification of all reported results and explicitly
  separates supported board evidence from non-claims such as production
  streaming/DMA throughput, cycle-level slowdown, and real-malware validation.
\end{itemize}

We focus on controlled RISC-V Linux workloads. The paper evaluates behavior
reconstruction and observer-boundary evidence; it does not claim malware
detection accuracy, uncontrolled real-malware validation, cycle-level
trace-on/trace-off slowdown, or production streaming transport.

Our artifact is available at \artifactroot.
